\scanpage{31}
$\overline{\mathcal V}={}$ catégorie des var. top. \inserted{à bord} avec immersions ouvertes

$\mathcal V^{\deleted{\unsure{*}}}={}$ catégorie des variétés topologiques sans bord

$\partial:\overline{\mathcal V}\longrightarrow\mathcal V$ foncteur bord topologique, $\deleted{\unsure{i|X}}X\longrightarrow\deleted{\unsure{i}}\operatorname{int}(X)\inserted{=X-\partial X}\unsure{=|:}\overline{\mathcal V}\longrightarrow\mathcal V$ \margin{foncteur ``pts intérieurs}

$\overline{\mathcal V}_{\leq n},\mathcal V_{\leq n},\overline{\mathcal V}_n,\mathcal V_n$ ss-catégories des variétés de dim $\leq n$ (resp. purement de dim $n$)

Une \uline{théorie de l'orientation} des variétés topologiques consiste en

(i) Un foncteur
\[
\Omega:\mathcal V^{\unsure{o|0|\circ}}\longrightarrow(\mathrm{Ens})
\]
\margin{On note $\Omega_n=\Omega|\mathcal V_n$\\ $\Omega(X)={}$ ensemble des orientations de $X$}

(ii) Une \deleted{\unsure{application|morphisme}} élément \deleted{de}
\[
\omega^+\in\Omega(e)\qquad(e\ \text{la variété ponctuelle})
\]

(iii) \deleted{\unsure{Pour tout|Donnée}} Un hom de foncteurs sur $\overline{\mathcal V}^{o}$
\[
\Omega(\operatorname{int}X)\longrightarrow\Omega(\partial X),
\]

Ces données satisfaisant aux conditions suivantes

\uline{1}$^\circ$) $\Omega$ est un \uline{faisceau} [i.e. $\forall X\in\operatorname{Ob}\mathcal V$, la restriction \inserted{$\Omega_X$} de $\Omega$ à la catégorie des ouverts de $X$ est un faisceau \unsure{$\}$} (faisceau d'orientation de $X$)]
% slanted
\margin{\unsure{nature locale|notion locale}}
\deleted{\unsure{de plus les} $\Omega_X\to X$ sont \unsure{couvrants|ouverts|surjectifs|revêtements} (\unsure{tte variété|les variétés} \unsure{a} \unsure{local$^{\rm mt}$ une orient.|localement une orientation})}

2$^\circ$) $\operatorname{card}(\Omega(e))=2$. [L'élément de $\Omega(e)$ distinct de $\omega^+$ est noté $\omega^-$. On définit une bijection $\Omega(e)\simeq\{+1,-1\}\subset\mathbb Z$ par $\omega^+\mapsto+1$, $\omega^-\mapsto-1$ : multiplicité (égale à $+1$ ou $-1$) d'un point orienté Catégorie $\vec{\mathcal V}^{\unsure{(\Omega)}}$ des variétés \unsure{top.|top} ``orientées'' (i.e. $\Omega$-orientées)]
% slanted
\margin{hyp. de \unsure{démarrage} des \unsure{récurrences|revêtements} (\unsure{normalis$^{\rm tion}$|normalisation})}

\deleted{\unsure{$3^\circ$})} Pour tte $\vec X\in\operatorname{Ob}\vec{\mathcal V}$ et $Y$ sous-variété de codim $1$, le faisceau $\unsure{\rho|p}_{Y/X}^{\unsure{(X'',\omega)|(X,\omega)}}$ des ``rives'' de $Y$ dans $X$ (qui est un rev. étale d'ordre $2$ de \unsure{$x$|$X$}) est envoyé, grâce à la donnée (iii), dans le faisceau d'orientation $\Omega_Y$. On trouve donc, \deleted{pour tout $X\in$} $\Omega(X)\to\operatorname{Hom}(\rho_{Y/X},\Omega_Y)$, et par faisceautisation
\[
\Omega_X|Y\longrightarrow\underline{\operatorname{Hom}}(\rho_{Y/X},\Omega_Y)
\]
ou encore
\[
\Omega_X\times\rho_{Y/X}\longrightarrow\Omega_Y
\]
% archive sheet number: 30
